\documentclass[10pt,a4paper]{amsart}
\usepackage[margin=19mm]{geometry}
\usepackage{amsmath,amssymb,mathtools}
\usepackage[scr=rsfs,cal=euler]{mathalfa}
\usepackage{xfrac}
\usepackage[unicode,hidelinks,bookmarks=false]{hyperref}
\newcommand{\ie}{i.e.\ }
\newcommand{\cf}{cf.\ }
\newcommand{\resp}{respectively\ }
\newcommand{\Subsection}{\subsection}
\providecommand{\nobreakdash}{\nobreak\hbox{-}}
\setlength{\emergencystretch}{2em}
\multlinegap0pt
\usepackage{amssymb}
\usepackage{enumerate}
\usepackage{tikz}
\newtheorem{proposition}{Proposition}
\title{On the generalized Tur\'an number of complete bipartite graphs}
\author{Oliver Janzer, Sean Longbrake, Liana Yepremyan}
\date{Source: arXiv:2606.09801v1 (CC BY 4.0)}
\begin{document}
\maketitle

\noindent\emph{Source and licence.} On the generalized Tur\'an number of complete bipartite graphs. Version arXiv:2606.09801v1 (CC BY 4.0), distributed by arXiv under the Creative Commons Attribution 4.0 International licence, http://creativecommons.org/licenses/by/4.0/, as recorded in the arXiv licence metadata for this exact version and shown on the abstract page https://arxiv.org/abs/2606.09801v1. Full licence notice: \texttt{licenses/arxiv-2606.09801-CC-BY-4.0.txt}.\par\smallskip

\noindent\textit{Editorial scope.} Counted unit: the article's proposition count (source0.tex lines 277--279). The article's complete original proof of it is delivered with it (source0.tex lines 363--368, one \texttt{proof} environment of 6 lines, which the article places away from the statement). No mathematics is written, repaired or re-derived: the statement and the proof are the article's own lines, in the article's own notation, and no retained line is substituted or re-flowed.  No further source-local context is needed: every cross-reference of every retained line resolves inside the two delivered spans.  Nothing was omitted from any retained span, so the package carries no omission record.  No retained line contains a citation, so the wrapper carries no bibliography and no bibliography entry.  No retained line contains a figure environment, an image-inclusion command or drawing code, so no image is delivered and the package declares no asset dependency.  Editorial formatting only, in addition: an AMS \texttt{amsart} preamble that restates the article's own declarations and macros used by the retained lines, namely the environments \texttt{proposition} on separate counters, together with the standard packages those lines need.  The retained environments print in this document as Proposition 1 in order of appearance, whereas the article prints the same items with the numbering of its own full document.  The complete native source of the article as distributed in the arXiv e-print for this exact version is bundled unchanged as \texttt{sources/202430/source0.tex}.
\begin{proposition} \label{prop:subspace count}
    The number of $k$-dimensional affine (projective) subspaces in affine (projective) $n$-space over $\mathbb{F}_q$ is $(1 +o(1))q^{(k + 1)(n - k)}$.
\end{proposition}

    \begin{proof}[Proof of (ii):]Let $H\in \bold{H}'$. Let $\bold{N}(H)$ be its neighborhood in $\bold{G}$. By Proposition~\ref{prop:subspace count}, $H$ contains $(1+o(1))q^6$ planes. Note that any three points in $\bold{N}(H)$ determine a plane contained in the hyperplane $H$, as $\bold{N}(H) \subseteq \bold{S}'\subseteq \varphi(S)$ contains no three collinear points. On the other hand, every plane contained in $H$ contains fewer than $t_0$ points from $\bold{N}(H)$ (as otherwise it would be dangerous for $\varphi(S)$ in which case we would have deleted $H$). Combining these inequalities, we have ${|\bold{N}(H)| \choose 3}/t_0^3\leq (1+o(1))q^6$, which implies  $|\bold{N}(H)|<8t_0q^2$.
    


Take now some point $P \in \bold{S}'$, and let $\bold{N}(P)$ be its neighborhood in $\bold{G}$. By Proposition~\ref{prop:subspace count}, $P$ is contained in $(1+o(1))q^6$ planes. Note that any three hyperplanes in $\bold{N}(P)$ intersect in a plane containing the point $P$, as $\bold{N}(P) \subseteq \bold{H}'\subseteq \bold{H}$ contains no three hyperplanes with a three-dimensional intersection. On the other hand, every plane containing $P$ is contained in fewer than $t_0$ hyperplanes from $\bold{N}(P)$ (as otherwise it would be dangerous for $\bold{H}$ in which case we would have deleted $P$). Combining these inequalities, we have ${|\bold{N}(P)| \choose 3}/t_0^3\leq (1+o(1))q^6$, which implies  $|\bold{N}(P)|<8t_0q^2$.
    \end{proof}

\end{document}
